\documentclass[letterpaper,12pt,oneside]{article}
\usepackage[top=1in,bottom=1in,left=1.25in,right=1.25in]{geometry}
\usepackage[T1]{fontenc}
\usepackage[utf8]{inputenc}
\usepackage[spanish,es-nodecimaldot,es-tabla]{babel}
\usepackage{lmodern,microtype,graphicx,caption,pdfpages}
\usepackage{amsmath,amssymb,booktabs,tabularx,longtable,xcolor,adjustbox}
\usepackage[hidelinks]{hyperref}
\numberwithin{figure}{section}
\widowpenalty=10000 \clubpenalty=10000
\setlength{\emergencystretch}{2em}
\setlength{\parskip}{0.3em}
\setlength{\textfloatsep}{12pt plus 3pt minus 3pt}
\setlength{\intextsep}{12pt plus 3pt minus 3pt}
\renewcommand{\topfraction}{.9}
\renewcommand{\bottomfraction}{.8}
\renewcommand{\textfraction}{.08}
\renewcommand{\floatpagefraction}{.75}
\setcounter{topnumber}{3} \setcounter{bottomnumber}{2} \setcounter{totalnumber}{5}
\raggedbottom
\captionsetup{font=small,labelfont=bf,skip=6pt,hypcap=false}
\newsavebox{\figbox}
\newdimen\figpageheight
\newdimen\normalpaperheight
\normalpaperheight=\paperheight
\makeatletter
\newcommand{\figura}[3]{%
 \sbox{\figbox}{\begin{minipage}{\linewidth}\centering
  \IfFileExists{#1.png}{%
   \adjustbox{max width=\linewidth}{\includegraphics{#1.png}}%
  }{\fbox{\parbox[c][.75in][c]{.90\linewidth}{\centering\small
   Evidencia pendiente\\[4pt]\texttt{\detokenize{#1.png}}}}}%
  \captionof{figure}{#2}\label{#3}
 \end{minipage}}%
 \ifdim\dimexpr\ht\figbox+\dp\figbox\relax>.85\textheight
  \par
  \ifx\@deferlist\@empty
   \ifx\@toplist\@empty
    \ifx\@botlist\@empty\else\clearpage\fi
   \else\clearpage\fi
  \else\clearpage\fi
  \figpageheight=\dimexpr\pagetotal+\ht\figbox+\dp\figbox+2in+24pt\relax
  \ifdim\figpageheight<\normalpaperheight\figpageheight=\normalpaperheight\fi
  \pdfpageheight=\figpageheight\paperheight=\figpageheight
  \textheight=\dimexpr\figpageheight-2in\relax
  \@colht=\textheight\@colroom=\textheight\vsize=\textheight\pagegoal=\textheight
  \noindent\usebox{\figbox}\par\clearpage
  \paperheight=\normalpaperheight\pdfpageheight=\normalpaperheight
  \restoregeometry
 \else
  \begin{figure}[htbp]\centering\usebox{\figbox}\end{figure}
 \fi}
\makeatother
\begin{document}
\includepdf[pages=1,pagecommand={}]{.build/caratula.pdf}
\pagenumbering{arabic}
\tableofcontents
\clearpage
\section{Introducción}
\subsection{Planteamiento del problema}
Se estudia la función de un destructor y su relación con el manejo de recursos. Además, se desarrolla una aplicación Java que recibe una oración e imprime cuántas palabras distintas aparecen más de una vez, sin distinguir mayúsculas de minúsculas ni considerar los signos de puntuación como parte de las palabras.
\subsection{Motivación}
El conteo de palabras permite identificar repeticiones sin revisar manualmente toda una oración. Para obtener resultados consistentes es necesario definir cómo se reconocen las palabras y separar esa tarea de la lectura e impresión. El estudio de los destructores permite distinguir la recuperación de memoria del cierre de otros recursos.
\subsection{Objetivos}
Se busca explicar qué hace un destructor, diferenciarlo de la recolección de basura en Java y resolver el conteo mediante objetos, ciclos, condicionales y un mapa de frecuencias. La aplicación debe mostrar el total de palabras duplicadas y las frecuencias correspondientes.

\section{Marco Teórico}
\subsection{Destructores y recursos}
En C++, un destructor es una función especial de una clase que se ejecuta cuando se destruye un objeto. Puede encargarse de liberar recursos adquiridos por ese objeto, como memoria reservada explícitamente. Para un objeto local de duración automática, la salida de su ámbito provoca su destrucción. El destructor no debe confundirse con cualquier método que borre datos: forma parte del ciclo de vida del objeto.\cite{destructor}

Java no ofrece destructores como los de C++. Su recolector puede recuperar la memoria de objetos que dejan de ser alcanzables; no se debe suponer que esto ocurre inmediatamente. El antiguo método \texttt{finalize()} no proporciona un cierre oportuno y está obsoleto, marcado para eliminación. Los recursos externos requieren mecanismos de cierre explícito; recuperar memoria no equivale a cerrar un archivo.\cite{finalizacion}

\subsection{Palabras y frecuencias}
En esta aplicación, una palabra es una secuencia de letras o dígitos. Los demás caracteres actúan como separadores; así, una coma entre dos palabras no las une. Se normalizan las letras a minúsculas y se conservan los acentos: \textit{Árbol} coincide con \textit{árbol}, pero \textit{si} y \textit{sí} permanecen distintos. Se emplean las operaciones de clasificación y conversión de \texttt{Character} para los caracteres habituales del español.\cite{caracter}

Un mapa asocia una clave con un valor. Aquí, la clave es la palabra normalizada y el valor es su número de apariciones. \texttt{HashMap} conserva claves únicas sin garantizar un orden de recorrido.\cite{mapa} Una palabra se considera duplicada cuando su frecuencia es mayor que uno; tres apariciones de la misma palabra aportan una unidad al total de palabras duplicadas.

\texttt{TreeSet} permite recorrer claves en su orden natural.\cite{conjunto} Este orden hace reproducible la presentación; no representa una ordenación lingüística especial para el español. \texttt{Map} y \texttt{Set} describen las operaciones de las colecciones utilizadas, sin requerir clases propias adicionales.

\section{Desarrollo}
\subsection{Investigación sobre destructores}
Un ejemplo conceptual es un objeto de C++ que reserva un bloque de memoria para almacenar texto: su destructor puede liberar ese bloque cuando termina la vida del objeto. Esto evita que el recurso quede ocupado después de dejar de utilizarlo.\cite{destructor} En la aplicación Java de esta práctica, el analizador únicamente conserva una oración y un mapa; no necesita un destructor ni un finalizador.\cite{finalizacion}

La síntesis anterior se apoya en documentación oficial. No se proporcionaron respuestas reales de las seis consultas, por lo que su comparación queda pendiente. Las figuras siguientes reservan las evidencias individuales. Añadir las capturas y recompilar permite mostrarlas, pero la comparación escrita requiere leer y contrastar sus respuestas.

La Figura~\ref{fig:consulta1} corresponde a la consulta de Luis.
\figura{Capturas/consulta1}{Consulta sobre destructores realizada por Luis}{fig:consulta1}

La Figura~\ref{fig:consulta2} corresponde a la consulta de Zoé.
\figura{Capturas/consulta2}{Consulta sobre destructores realizada por Zoé}{fig:consulta2}

La Figura~\ref{fig:consulta3} corresponde a la consulta de Alex.
\figura{Capturas/consulta3}{Consulta sobre destructores realizada por Alex}{fig:consulta3}

La Figura~\ref{fig:consulta4} corresponde a la consulta de Joseph.
\figura{Capturas/consulta4}{Consulta sobre destructores realizada por Joseph}{fig:consulta4}

La Figura~\ref{fig:consulta5} corresponde a la consulta de Hanniel.
\figura{Capturas/consulta5}{Consulta sobre destructores realizada por Hanniel}{fig:consulta5}

La Figura~\ref{fig:consulta6} corresponde a la consulta de Alan.
\figura{Capturas/consulta6}{Consulta sobre destructores realizada por Alan}{fig:consulta6}

\subsection{Organización de la aplicación}
\texttt{Main} lee una línea con \texttt{Scanner}, crea el analizador, solicita el conteo e imprime el total. Después pide mostrar las palabras duplicadas con ordenación activada. Su dependencia de \texttt{AnalizadorPalabras} se limita al objeto local que utiliza; no declara atributos ni constructor explícito.

La Figura~\ref{fig:flujo1} presenta la secuencia de la aplicación.
\figura{Diagramas/flujo1}{Flujo de la clase Main}{fig:flujo1}
La Figura~\ref{fig:clases1} presenta la tabla de Main; la fila de atributos está vacía porque no existen atributos declarados.
\figura{Diagramas/clases1}{Tabla de la clase Main}{fig:clases1}

\texttt{AnalizadorPalabras} conserva los atributos públicos \texttt{oracion} y \texttt{frecuencias}, siguiendo la visibilidad sugerida en el enunciado. Su constructor guarda la oración recibida e inicializa el mapa vacío. \texttt{contarPalabras} vacía el mapa antes de comenzar, de modo que repetir el conteo no acumule resultados de una ejecución anterior.

\subsection{Algoritmo de conteo}
Se recorre la oración carácter por carácter. Cuando se encuentra una letra o un dígito, se agrega su versión en minúscula a una palabra temporal. Cuando aparece un separador, se registra la palabra únicamente si no está vacía: si su clave existe, se incrementa la frecuencia; si no existe, se guarda con frecuencia uno. Después se vacía la palabra temporal para comenzar la siguiente.

El recorrido incluye una posición adicional al final. En esa posición no se lee ningún carácter: se ejecuta el mismo cierre usado para un separador. Así se registra la última palabra aunque la oración termine directamente en una letra. Los separadores consecutivos no generan claves vacías. Una línea vacía o formada solo por signos deja el mapa vacío y produce cero duplicadas, sin solicitar reintentos que el enunciado no exige.

Durante el recorrido, el mapa contiene las frecuencias de las palabras completas ya cerradas y la variable temporal contiene únicamente el fragmento de la palabra actual. Esta relación explica por qué cada palabra se registra una sola vez al encontrar su límite.

La Figura~\ref{fig:flujo2} resume el conteo. Los pasos posteriores al nodo de recorrido describen sus operaciones internas condicionales: acumular y registrar corresponden a casos distintos, no a acciones que se ejecuten juntas para cada carácter. El cierre incluye la posición final y solo registra palabras no vacías.
\figura{Diagramas/flujo2}{Conteo de frecuencias en AnalizadorPalabras}{fig:flujo2}

\texttt{obtenerNumeroDuplicadas} recorre los valores del mapa y suma una unidad por cada frecuencia mayor que uno. \texttt{mostrarDuplicadas} recorre las claves e imprime únicamente las que cumplen esa condición. Si recibe verdadero, utiliza un \texttt{TreeSet} local; si recibe falso, utiliza el recorrido del mapa sin prometer un orden. Consultar o mostrar el resultado no modifica las frecuencias. Estos métodos son operaciones independientes que Main invoca después del conteo.

La Figura~\ref{fig:clases2} muestra los atributos y las operaciones reales del analizador.
\figura{Diagramas/clases2}{Tabla de la clase AnalizadorPalabras}{fig:clases2}

\section{Resultados}
\subsection{Oración con repeticiones}
La primera entrada es \texttt{¡Hola,hola! HOLA; Java, java. Árbol árbol.} Las frecuencias calculadas son tres para \textit{hola}, dos para \textit{java} y dos para \textit{árbol}. Por tanto, el total esperado es tres palabras duplicadas. El número de apariciones adicionales sería cuatro, pero esa es una medida distinta de la solicitada aquí. La prueba reúne cambios de mayúsculas, acentos, puntuación y una coma sin espacios.

La salida imprime el total tres y después las líneas \texttt{hola: 3}, \texttt{java: 2} y \texttt{árbol: 2}, en ese orden. La Figura~\ref{fig:ejecucion1} está destinada a la captura real de esta ejecución.
\figura{Capturas/ejecucion1}{Ejecución con tres palabras distintas duplicadas}{fig:ejecucion1}

\subsection{Oración sin repeticiones}
La segunda entrada es \texttt{Sol, luna y mar.} Cada palabra aparece una vez: se imprime cero y no se agregan líneas de palabras después del encabezado de frecuencias. Esta prueba contrasta la existencia y la ausencia de duplicadas. La Figura~\ref{fig:ejecucion2} está destinada a la captura real correspondiente.
\figura{Capturas/ejecucion2}{Ejecución sin palabras duplicadas}{fig:ejecucion2}

Las dos pruebas se verificaron mediante compilación y ejecución Java, comparando la salida completa con los resultados esperados, incluidos los mensajes de lectura y los saltos de línea. Las capturas coloreadas corresponden al proceso posterior del motor del usuario; no se sustituyen por imágenes fabricadas. No se dispuso de main.py para validar o renderizar su geometría.

\section{Conclusiones}
La normalización determina qué apariciones pertenecen a una misma palabra. Unificar mayúsculas y utilizar los signos como separadores evita dividir artificialmente sus frecuencias o unir palabras vecinas. El mapa permite distinguir entre el número de palabras diferentes que se repiten y el número de repeticiones adicionales; las pruebas reflejan esa diferencia y el caso sin duplicadas.

Separar la interacción de Main del comportamiento del analizador mantiene una correspondencia clara entre datos y operaciones. Por otra parte, el estudio de destructores muestra que la liberación de recursos depende del lenguaje: en Java, la recolección de memoria no justifica depender de un finalizador para cerrar recursos externos. La comparación de las consultas individuales sigue pendiente de contar con sus respuestas reales.

\begin{thebibliography}{9}
\bibitem{destructor} Microsoft. \textit{Destructors (C++)}. Microsoft Learn. \url{https://learn.microsoft.com/en-us/cpp/cpp/destructors-cpp}.
\bibitem{finalizacion} OpenJDK. \textit{JEP 421: Deprecate Finalization for Removal}. \url{https://openjdk.org/jeps/421}.
\bibitem{caracter} Oracle. \textit{Character}. Java SE 21 API. \url{https://docs.oracle.com/en/java/javase/21/docs/api/java.base/java/lang/Character.html}.
\bibitem{mapa} Oracle. \textit{HashMap}. Java SE 21 API. \url{https://docs.oracle.com/en/java/javase/21/docs/api/java.base/java/util/HashMap.html}.
\bibitem{conjunto} Oracle. \textit{TreeSet}. Java SE 21 API. \url{https://docs.oracle.com/en/java/javase/21/docs/api/java.base/java/util/TreeSet.html}.
\end{thebibliography}

\end{document}
